% TOP_PROBLEM: {"id": 30006770, "title": "Smooth anisotropic Calderón uniqueness from full boundary data", "category_id": 9, "category": {"id": 9, "name": "pde", "display_name": "Partial Differential Equations", "description": "PDEs and their applications in physics and geometry.", "slug": "pde", "order_index": 9, "created_at": "2026-07-31T15:26:25.671Z"}, "status": "open"}
% FIRST_POSTED: 2026-09-27
% !TeX program = lualatex
% Compile twice: lualatex 30006770.tex
% All references and prose are embedded; no BibTeX database or external inputs.
\documentclass[11pt,a4paper]{article}
\usepackage[margin=24mm,headheight=14pt]{geometry}
\usepackage{fontspec}
\setmainfont{Latin Modern Roman}
\setsansfont{Latin Modern Sans}
\setmonofont{Latin Modern Mono}
\usepackage{amsmath,amssymb,mathtools}
\usepackage{unicode-math}
\setmathfont{Latin Modern Math}
\usepackage{microtype}
\usepackage{enumitem,xurl,needspace}
\usepackage[dvipsnames]{xcolor}
\usepackage{hyperref}
\hypersetup{unicode=true,colorlinks=true,linkcolor=MidnightBlue,urlcolor=MidnightBlue,
 pdftitle={500 Mathematical Problem Statements},pdfauthor={Source-annotated compilation},bookmarksnumbered=true}
\usepackage{bookmark}
\usepackage{tocloft}
\setlength{\cftsecnumwidth}{3em}
\cftsetpnumwidth{2.5em}
\cftsetrmarg{4em}
\usepackage{titlesec}
\titleformat{\section}{\Large\bfseries\raggedright}{\thesection}{0.7em}{}
\usepackage{fancyhdr}
\pagestyle{fancy}\fancyhf{}
\fancyhead[L]{\small\sffamily Mathematical problem statements}
\fancyhead[R]{\small Source edition 22}
\fancyfoot[C]{\thepage}
\setlength{\parindent}{0pt}
\setlength{\parskip}{5pt plus 1pt}
\setlength{\emergencystretch}{3em}
\setcounter{tocdepth}{1}
\setcounter{secnumdepth}{1}
\setlist[itemize]{leftmargin=3em,itemsep=5pt,topsep=4pt,parsep=0pt}
\allowdisplaybreaks
\newcommand{\N}{\mathbb{N}}
\newcommand{\Z}{\mathbb{Z}}
\newcommand{\Q}{\mathbb{Q}}
\newcommand{\R}{\mathbb{R}}
\newcommand{\C}{\mathbb{C}}
\newcommand{\F}{\mathbb{F}}
\newcommand{\A}{\mathbb{A}}
\newcommand{\PP}{\mathbb P}
\newcommand{\E}{\mathbb{E}}
\newcommand{\eps}{\varepsilon}
\newcommand{\dd}{\,\mathrm d}
\newcommand{\sref}[2]{\hyperlink{source:#1:#2}{[S#2]}}
\newcommand{\eref}[2]{\hyperlink{extra:#1:#2}{[E#2]}}
% Turn the next switch off for a shorter reading copy. The quotations preserve
% every exactTarget field, including qualifications omitted from a short summary.
\newif\ifshowcatalogue
\showcataloguetrue
\newcommand{\cataloguescope}[1]{\ifshowcatalogue
 \paragraph{Exact scope in the supplied catalogue.}
 \begin{quote}\small #1\end{quote}\fi}

% Standalone notebook wrapper; original problem section retained verbatim.
\numberwithin{equation}{section}
\hypersetup{pdftitle={Research attempt -- problem 30006770},
 pdfauthor={Research notebook; source compilation retained}}
\fancyhead[L]{\small\sffamily Research notebook}
\fancyhead[R]{\small\sffamily Problem 30006770 | 27 September 2026}
\begin{document}
\setcounter{section}{0}
% ================================================================
% problemId: 30006770
% Canonical problem ID: 30006770
% exactTarget SHA-256: 5cfb58563fd37180b8527904e281cfa5f02bb3385d78adf9bc1cf4743f3c2bbe
\clearpage
\section*{\texorpdfstring{Smooth anisotropic Calderón uniqueness from full boundary data}{Smooth anisotropic Calderón uniqueness from full boundary data}}\label{30006770}
\subsection{Definitions and mathematical statement}
For a smooth metric $g$ on a compact connected manifold $M$ with boundary, solve the Dirichlet problem $\Delta_g u=0$, $u|_{\partial M}=f$, and define the Dirichlet-to-Neumann operator
$\Lambda_g f=\partial_{\nu_g}u|_{\partial M}$, with outward metric-unit normal $\nu_g$. In dimension $n\ge3$, the target is
\[
 \Lambda_{g_1}=\Lambda_{g_2}\quad\Longrightarrow\quad
 g_2=\Phi^*g_1\text{ for some }\Phi|_{\partial M}=\mathrm{id}.
\]
Both metrics live on the same smooth manifold and the operators are compared on the entire boundary using its identity marking. Pullback by a boundary-fixing diffeomorphism is the unavoidable gauge ambiguity. No real-analyticity, product structure, or partial-boundary restriction is included.
\par\smallskip\textit{Formulation and concept references:} \sref{30006770}{1}.
\cataloguescope{For every compact connected smooth manifold M of dimension n>=3 with smooth nonempty boundary and every pair of smooth Riemannian metrics g1,g2 on that same M, let Lambda\_gj f be the outward metric-unit normal derivative of the unique solution of Delta\_gj u=0 with u|boundary=f. If Lambda\_g1=Lambda\_g2 as operators C-infinity(boundary M) to C-infinity(boundary M), compared under the identity marking on the whole boundary, must there be a smooth diffeomorphism Phi:M->M fixing the boundary pointwise with g2=Phi\textasciicircum{}*g1?}
\subsection{Short English statement}
Do all boundary voltage-to-flux measurements determine a smooth interior metric, apart from a change of coordinates that fixes the boundary?
% BEGIN RESEARCH ADDITION ATTEMPT-176
\subsection{Research attempt}

\paragraph{Current frontier.}
The 2026 formulation paper explicitly distinguishes the general smooth
problem from its quasianalytic, symmetric, and ordered settings
\sref{30006770}{1}, Introduction and Section 1.1. Its new results are manuscript
results with additional hypotheses, not a proof of the present universal
assertion. Boundary determination fixes jets in boundary-normal coordinates;
it does not identify arbitrary smooth metrics on an interior collar merely
from equality of those jets. We use the standard energy and unique
continuation mechanisms discussed in that source, but no unverified global
rigidity theorem from it.

\paragraph{Attack route.}
A boundary-jet route requires a continuation principle unavailable for
arbitrary smooth coefficients. A monotonicity route works if the metric
comparison produces a sign-definite energy defect. A linearized route can
use products of harmonic functions, but must distinguish injectivity from
an invertible nonlinear reconstruction. We pursue the latter two routes:
prove an ordered conformal case directly, compute the Euclidean linearized
kernel in a gauge, and test the attempted passage to the full problem.

\paragraph{Attempt: an ordered conformal case on the stated manifold.}
Suppose, as additional hypotheses, that
$g_2=e^{2\varphi}g_1$, $\varphi\ge0$, and $\varphi$ vanishes on a
neighborhood of the boundary. Put $a=e^{(n-2)\varphi}$.
The Dirichlet energy of $g_2$, written using $g_1$, is
\[
 E_2(u)=\int_M a|\nabla u|_{g_1}^2\,dV_{g_1},
 \qquad E_1(u)=\int_M|\nabla u|_{g_1}^2\,dV_{g_1}.
\]
For boundary data $f$, let $u_j$ minimize its respective energy. The metrics
coincide near the boundary, so the surface measures and outward normals
agree there. Equality of the \emph{normal-derivative} DN maps in the question
therefore gives equality of the boundary energy pairings, and hence
$E_2(u_2)=E_1(u_1)$. Minimality of $u_1$ yields
\begin{equation}\label{eq176-monotone}
 0\le\int_M(a-1)|\nabla u_2|_{g_1}^2\,dV_{g_1}
     =E_2(u_2)-E_1(u_2)
     \le E_2(u_2)-E_1(u_1)=0.
\end{equation}
If $\varphi$ is positive somewhere, it is bounded below by a positive
number on an interior ball. Thus $u_2$ is constant on that ball. Interior
unique continuation for smooth uniformly elliptic scalar equations then
makes $u_2$ constant on connected $M$. Choose nonconstant smooth boundary
data to get a contradiction. The unique continuation input, including its
applicability to smooth divergence-form coefficients, is the one recalled
in \sref{30006770}{1}, Section 5.1. Hence $\varphi=0$.
The case $\varphi\le0$ follows by interchanging the metrics.

This proves a genuine special-case uniqueness statement using the exact
boundary convention. It does not assume that all tensor-valued metric
differences can be made conformal or ordered by a boundary-fixing
coordinate change. Removing that additional structure removes the sign
in \eqref{eq176-monotone}.

\paragraph{Attempt: the Euclidean linearized gauge calculation.}
Let $\Omega\subset\R^n$, $n\ge3$, be a smooth bounded connected domain,
and let $h\in C_c^\infty(\Omega;\operatorname{Sym}_n)$.
For $g_\eps=I+\eps h$, the conductivity density is
\begin{equation}\label{eq176-density}
 \sqrt{\det g_\eps}\,g_\eps^{-1}
   =I+\eps B+O(\eps^2),\qquad
 B=\tfrac12(\operatorname{tr}h)I-h.
\end{equation}
For harmonic $u,v$, differentiation of the energy pairing gives the
linearized DN form
\[
                       \int_\Omega B\nabla u\cdot\nabla v\,dx.
\]
The derivatives of the harmonic extensions have zero boundary trace and
pair to zero with harmonic functions by integration by parts. Compact
support of $h$ also prevents a derivative of the boundary surface measure.
Thus vanishing of the linearized data implies that this integral is zero
for every pair of smooth harmonic functions.

Extend $B$ by zero and use complex harmonic exponentials
$u=e^{\zeta\cdot x}$, $v=e^{\eta\cdot x}$ with
$\zeta\cdot\zeta=\eta\cdot\eta=0$. Complex functions are legitimate by
complex-linear extension of the real bilinear identity. Fix a nonzero real
frequency $\xi$, choose a real $w\perp\xi$ of length $|\xi|/2$, and set
\[
 \zeta=i\xi/2+w,\qquad \eta=i\xi/2-w.
\]
They are null vectors and sum to $i\xi$. Symmetry of $B$ cancels the cross
terms, so, writing $\widehat B(\xi)=\int e^{i\xi\cdot x}B(x)\,dx$,
\begin{equation}\label{eq176-fourier}
      \tfrac14\xi^t\widehat B(\xi)\xi+
                           w^t\widehat B(\xi)w=0.
\end{equation}
Rotate coordinates so $\xi=|\xi|e_1$. Since \eqref{eq176-fourier} holds
for every $w$ on the transverse sphere, polarization gives
\[
 \widehat B_{ab}=-\widehat B_{11}\delta_{ab}
                       \qquad(2\le a,b\le n).
\]
The invertible linear relation in \eqref{eq176-density} is
\[
 \widehat h=\frac{\operatorname{tr}\widehat B}{n-2}I-\widehat B.
\]
Because $\operatorname{tr}\widehat B=-(n-2)\widehat B_{11}$, it follows
that $\widehat h_{ab}=0$ for all transverse indices $a,b\ge2$.

Now impose the explicit linear gauge $\operatorname{div}h=0$.
It implies $\xi_i\widehat h_{ij}=0$, so in these coordinates all first-row
and first-column entries vanish too. Hence $\widehat h(\xi)=0$ for every
$\xi\ne0$, and continuity gives the same at zero. Fourier injectivity
proves $h=0$. This is a complete injectivity calculation for compactly
supported, divergence-free infinitesimal perturbations, not the nonlinear
problem on an arbitrary manifold.

The omitted linear directions are exactly of the expected potential form
at each nonzero frequency: a symmetric matrix with only first row and
column nonzero can be written as $\xi\otimes v+v\otimes\xi$.
Indeed a boundary-fixing infinitesimal diffeomorphism produces
$h=\nabla V+(\nabla V)^t$, which has zero linearized DN data. This also
follows directly by differentiating the exact pullback invariance. The
calculation does not claim a global gauge-fixing theorem with estimates
from the frequency description alone.

\paragraph{Attempt: trying to close the nonlinear gap.}
A possible local strategy would choose a gauge and prove an estimate
controlling $h$ by the linearized data in spaces where the nonlinear
remainder is smaller than that control. The injectivity above supplies no
bounded inverse in such spaces. Harmonic exponentials grow with frequency;
using them in estimates incurs costs absent from a bare zero-data argument.
The assertion that an injective derivative automatically gives a local
inverse in an infinite-dimensional problem is false without additional
range and stability hypotheses.

The boundary-jet route has its own explicit failure. A smooth function
$e^{-1/t^2}$ for $t>0$, extended by zero for $t\le0$, has every derivative
zero at zero but is not zero on any right neighborhood. Multiplying this
by a tensor in a collar produces smooth perturbations invisible to their
boundary Taylor series, while not proving anything about their full DN
maps. The example refutes inference from equal jets to equal smooth metrics,
not uniqueness from the full boundary data.

\paragraph{Stress test.}
Dimension two is an important negative check on the algebra:
$n-2=0$ makes the tensor inversion in \eqref{eq176-density} impossible,
and conformal rescaling with factor one at the boundary preserves the
harmonic equation and its measured normal derivative. That mechanism is
excluded by the stated $n\ge3$ hypothesis; it is not a counterexample here.

The weak boundary flux is a density whereas the catalogue DN map is a
function. In both calculations the metrics agree near the boundary, so
the distinction has been resolved explicitly rather than ignored. General
equality of DN operators would first need the appropriate boundary
identification, as discussed in \sref{30006770}{1}. Neither positivity of a
conformal factor nor a Euclidean Fourier representation is available for
two unrestricted smooth metrics in the target.

\paragraph{Outcome.}
\textbf{Outcome: Exploratory approach with unresolved gap.}

The attempt supplies proofs of an ordered conformal special case and
linearized Euclidean injectivity in a divergence-free compact-support
gauge. These are diagnostic derivations of familiar mechanisms, not claims
of new general inverse-problem theorems. The unresolved steps are obtaining
sufficient harmonic probing and quantitative nonlinear control for arbitrary
smooth geometry, or another global rigidity mechanism without order or
quasianalyticity. No full uniqueness theorem or counterexample is obtained.
% END RESEARCH ADDITION ATTEMPT-176
\subsection{Sources}
\begin{itemize}\raggedright
\item[\textnormal{[S1]}]\hypertarget{source:30006770:1}{} Quasianalyticity and geometric rigidity in anisotropic Calderon's problem; arXiv2608.15813v1 August16,2026.
\url{https://arxiv.org/pdf/2608.15813v1}.
{\small\textit{Catalogue formulation source.}}
\end{itemize}

\end{document}
